\section{Pillowcase Floer homology}\label{sec:hhk}

This section collects the parts of Hedden, Herald and Kirk's Lagrangian Floer theory in the
pillowcase~\cite{HHK2} that we use, together with Daemi and Scaduto's S-complexes~\cite{DSEquiv}.
Results of~\cite{HHK2} are cited by their numbering in arXiv:1501.00028v1. Coefficients in Floer
complexes are in $\F_2$ throughout.

\subsection{The pillowcase}\label{ss:pillowcase}
Let $G\cong\Z^2\rtimes\Z/2$ be the group of isometries of $\R^2$ generated by
$(\gamma,\theta)\mapsto(\gamma+2\pi,\theta)$, $(\gamma,\theta)\mapsto(\gamma,\theta+2\pi)$ and
$(\gamma,\theta)\mapsto(-\gamma,-\theta)$. The \emph{pillowcase} is the quotient $P=\R^2/G$, a
topological $2$-sphere; the quotient map $\R^2\to P$ is a branched covering, branched over the images of
the lattice points $(\pi\Z)^2$, which are the four \emph{corners} of $P$~\cite[\S3.1]{HHK2}. The rectangle
$[0,\pi]\times[0,2\pi]$ is a fundamental domain; we call $[0,\pi]\times[0,\pi]$ the front face,
$[0,\pi]\times[\pi,2\pi]$ the back face, and the images of the lines $\gamma\in\pi\Z$ the left and right
edges. Put $\Pst=(\R^2\setminus(\pi\Z)^2)/G$, oriented by $d\gamma\wedge d\theta$. The restriction
\[
\R^2\setminus(\pi\Z)^2\longrightarrow\Pst
\]
is a regular covering with deck group $G$. Paths, loops and immersed disks in $\Pst$ therefore lift to
$\R^2\setminus(\pi\Z)^2$, and we shall work with such lifts throughout.

The pillowcase is the traceless character variety of the four-punctured sphere: the map
$\psi(\gamma,\theta)\colon a\mapsto\mathbf i$, $b\mapsto e^{\gamma\mathbf k}\mathbf i$, $c\mapsto e^{\theta\mathbf
k}\mathbf i$, $d\mapsto e^{(\theta-\gamma)\mathbf k}\mathbf i$ induces a homeomorphism of $P$ with
$R(S^2,\{a,b,c,d\})$, the space of conjugacy classes of $\SU$ representations of
$\pi_1(S^2\setminus\{a,b,c,d\})=\langle a,b,c,d\mid ba=cd\rangle$ sending each of $a,b,c,d$ to a traceless
element, and the corners correspond to the abelian representations~\cite[Prop.~6.1]{HHK2},
\cite[Prop.~3.1]{HHK1}.

On $\R^2$ we use the function
\[
\varphi=\theta-\gamma .
\]
Translations in $G$ change $\varphi$ by elements of $2\pi\Z$, and $(\gamma,\theta)\mapsto(-\gamma,-\theta)$ changes
its sign. The \emph{diagonal arc} $\Delta=\{(\gamma,\gamma)\mid\gamma\in[0,\pi]\}\subset P$~\cite[(12)]{HHK2}
joins the corners $(0,0)$ and $(\pi,\pi)$, and its preimage in $\R^2$ is the union of the lines
\[
\Delta_k=\{\varphi=2\pi k\},\qquad k\in\Z .
\]
The lines $\{\varphi=\pi+2\pi k\}$ are lifts of the arc of slope one on the back face, which is not part of
$\Delta$. Every lattice point lies on a line $\{\varphi\in\pi\Z\}$; the lattice points on $\Delta_k$ are the
points $(j\pi,j\pi+2\pi k)$.

\subsection{The earring curve}\label{ss:earring}
For $\epsilon>0$ let $L_0=L_0^{\epsilon,0}\colon S^1\to\Pst$ be the immersion defined by
\begin{equation}\label{eq:L0}
\tilde L_0(t)=\bigl(t+\epsilon\sin t+\tfrac\pi2,\ t-\epsilon\sin t+\tfrac\pi2\bigr),\qquad t\in[0,2\pi],
\end{equation}
followed by the branched covering~\cite[Def.~3.4, (13)]{HHK2}; as in~\cite[\S11]{HHK2} we take the perturbation
function $g$ of~\cite[(13)]{HHK2} to be zero. It is the restriction to the boundary of the
traceless character variety of the trivial tangle with an earring, after a small holonomy
perturbation~\cite[Thm.~6.2]{HHK2}. For small $\epsilon$ it is self-transverse with one double point, at
$(\pi/2,\pi/2)$; it lies in an $O(\epsilon)$-neighbourhood of $\Delta$ and passes around the two corners
$(0,0)$ and $(\pi,\pi)$ through the back face. It is an unobstructed circle with
$\mu(L_0,\ell_1)=0=z(L_0)$ in the notation of \S\ref{ss:restricted}~\cite[pp.~16,~18]{HHK2}.

Since $\tilde L_0(t+2\pi)=\tilde L_0(t)+(2\pi,2\pi)$, the preimage of $L_0$ in $\R^2$ is the union of the
curves
\[
S^k_+=\tilde L_0+(0,2\pi k),\qquad S^k_-=-\tilde L_0+(0,2\pi k),\qquad k\in\Z,
\]
which we call the \emph{strands} of $L_0$; each is a properly embedded line lying within $\sqrt2\,\epsilon$ of
$\Delta_k$, and a translation by $(2\pi a,2\pi b)$ carries $S^k_\pm$ to $S^{k+b-a}_\pm$. Reparametrizing,
$S^0_-=\{(t+\frac\pi2-\epsilon\sin t,\ t+\frac\pi2+\epsilon\sin t)\}$. The slope of $S_+$ at the point with
parameter $t$ is $(1-\epsilon\cos t)/(1+\epsilon\cos t)$, and that of $S_-$ is its reciprocal. Hence on the
front strip $0<\gamma<\pi$ the strand $S^k_+$ has slope less than $1$ and $S^k_-$ has slope greater than $1$;
in $P$ these are the two branches of $L_0$ along the interior of $\Delta$~\cite[\S5.1]{HHK2}. The two
strands near $\Delta_k$ cross each other exactly at the points $(\frac\pi2+j\pi,\frac\pi2+j\pi+2\pi k)$, and they
pass each lattice point of $\Delta_k$ on opposite sides at distance $\sqrt2\,\epsilon$: for instance $S^0_+$ passes
$(0,0)$ through $(-\epsilon,\epsilon)$ and $S^0_-$ through $(\epsilon,-\epsilon)$, and the sides alternate from one
lattice point of $\Delta_k$ to the next.

\subsection{Restricted curves and admissible pairs}\label{ss:restricted}
Let $\ell_1$ be the line field of slope one on $\R^2$; it descends to $\Pst$. For a path $\alpha$ immersed in
$\Pst$, the Maslov index $\mu(\alpha,\ell_1)$ is the signed number of times the tangent line of $\alpha$ passes
$\ell_1$, counterclockwise passages counting $+1$, with the half-open endpoint convention
of~\cite[\S2.3]{HHK2}: if $\vartheta(t)$ is a continuous choice of the angle of the tangent line minus $\pi/4$,
then $\mu=\lfloor(\vartheta(1)-\varsigma)/\pi\rfloor-\lfloor(\vartheta(0)-\varsigma)/\pi\rfloor$ for small $\varsigma>0$.
For two lines $\ell_0,\ell_1'$ and a line $\ell$, all pairwise transverse, the \emph{triple index}
$\tau(\ell_0,\ell_1',\ell)\in\{0,1\}$ is $1$ exactly when the clockwise rotation of $\ell_0$ to $\ell_1'$ passes
through $\ell$~\cite[Def.~2.5]{HHK2}. Let $z\in H^1(\Pst;\Z/4)$ be the class assigning $1$ to a small loop
around a corner, oriented counterclockwise~\cite[Def.~3.1]{HHK2}. A loop winding once around the
pillowcase in the open strip between the two edges encircles two corners, so $z$ takes the value $2$ on
it.

An immersed circle in $\Pst$ is \emph{unobstructed} if every lift of it to the universal cover of $\Pst$ is a
properly embedded line; by a lemma of Abouzaid~\cite[Lemma~2.2]{Abouzaid} this holds if and only if it
is homotopically essential and contains no \emph{fishtail}, that is, no subinterval whose endpoints map to
the same point and which maps to a nullhomotopic loop~\cite[Def.~2.1]{HHK2}. A \emph{proper immersion}
of an interval is the image of a smooth immersion $I\to\R^2$ sending the endpoints to lattice points and
the interior into $\R^2\setminus(\pi\Z)^2$; its slopes at the endpoints are its \emph{limiting
slopes}~\cite[Def.~3.3]{HHK2}.

\begin{definition}[{\cite[Def.~3.6 and p.~33]{HHK2}}]\label{def:restricted}
A \emph{restricted immersed circle} is an unobstructed immersed circle $L\colon S^1\to\Pst$ with
$\mu(L,\ell_1)+z(L)\equiv0\pmod4$. A \emph{restricted immersed arc} is a proper immersion of an interval
whose restriction to the complement of small discs around the corners contains no fishtail. A
\emph{restricted immersed $1$-manifold} is a disjoint union of one restricted immersed arc and finitely
many restricted immersed circles.
\end{definition}

The image of a straight segment joining two lattice points whose interior misses the lattice is a
restricted immersed arc~\cite[p.~18]{HHK2}.

\begin{definition}[{\cite[Def.~3.11]{HHK2}}]\label{def:admissible}
A pair $(L_0,L_1)$ of restricted immersed curves is \emph{admissible} if at least one of them is a circle,
if $L_0$ and $L_1$ intersect transversely, and if, whenever loops $\alpha_0$ in the domain of $L_0$ and
$\alpha_1$ in the domain of $L_1$ have $L_0\circ\alpha_0$ freely homotopic to $L_1\circ\alpha_1$ in $\Pst$, both
$\alpha_0$ and $\alpha_1$ are nullhomotopic.
\end{definition}

\subsection{The complex and its grading}\label{ss:complex}
For an admissible pair, $C(L_0,L_1)$ is the $\F_2$-vector space on $L_0\cap L_1$. A \emph{bigon} from $p$ to
$q$ is an orientation-preserving immersion of a disc into $\Pst$ with two convex corners, at $p$ and $q$,
whose boundary, traversed counterclockwise, runs along $L_0$ from $p$ to $q$ and then along $L_1$ from $q$
back to $p$; equivalent bigons are identified~\cite[Def.~2.10, (5)]{HHK2}. The differential is
$\partial p=\sum_q\#\mathcal M(p,q)\,q$, the count of bigons from $p$ to $q$ modulo $2$. By~\cite[Cor.~3.13]{HHK2}
each $\mathcal M(p,q)$ contains at most one bigon. A bigon lifts to an immersed disc in $\R^2\setminus(\pi\Z)^2$,
so its boundary lies on one component of $L_1$ and the two generators lie on the same component; in
particular there is no differential between generators on different components of
$L_1$~\cite[p.~37]{HHK2}. We write $HF(L_0,L_1)$ for the homology.

The complex carries a relative $\Z/4$-grading $\gr(p,q)=\gr(p)-\gr(q)$, defined for generators on the same
component of $L_1$~\cite[Def.~3.7]{HHK2}; the differential lowers it by one, and a bigon from $p$ to $q$
forces $\gr(p,q)=1$~\cite[Prop.~3.8]{HHK2}. It is computed as follows~\cite[Lemma~3.10]{HHK2}: if
$\alpha_0$ is a path in $L_0$ from $p$ to $q$ and $\alpha_1$ a path in $L_1$ from $q$ to $p$, then
\begin{equation}\label{eq:grading15}
\gr(p,q)=\mu(\alpha_0,\ell_1)+\mu(\alpha_1,\ell_1)+\tau(T_pL_0,T_pL_1,\ell_1)-\tau(T_qL_0,T_qL_1,\ell_1)
+z(\alpha_0*\alpha_1)\pmod4 .
\end{equation}
Two consequences are recorded in~\cite[\S\S5.1--5.2]{HHK2}.
\begin{enumerate}
\item If $L_1$ meets the interior of $\Delta$ transversely at $x$, then for small $\epsilon$ there are exactly two
generators near $x$: $x^+$ on the branch of $L_0$ of slope less than $1$ and $x^-$ on the branch of slope
greater than $1$. They satisfy $\gr(x^+,x^-)=1$.
\item If $p_+,q_+$ lie on the branch of slope less than $1$ on the front face, $\alpha_0$ runs along that branch,
and $L_1$ is transverse to $\Delta$ at both points, then
\begin{equation}\label{eq:grading16}
\gr(p_+,q_+)=\mu(\alpha_1,\ell_1)+z(\alpha_0*\alpha_1)\pmod4 ,
\end{equation}
and the same formula holds for two points on the branch of slope greater than $1$.
\end{enumerate}
In terms of lifts, $x^\pm$ are the intersections of a lift of $L_1$ with the strands $S^k_\pm$ near a point
of $\Delta_k$ on the front strip.

Suppose $L_1$ is a restricted immersed arc with an endpoint at the corner $(0,0)$ whose limiting slope is
bounded away from $1$. Then for small $\epsilon$ there is a unique intersection point $r_+$ of $L_0$ and $L_1$
converging to that corner as $\epsilon\to0$~\cite[Lemma~5.1]{HHK2}. When $L_1$ comes from a knot $K$ as in
\S\ref{ss:decomp}, the grading of the component containing the arc is made absolute by
\begin{equation}\label{eq:grading19}
\gr(r_+)=\sigma(K)\pmod4 ,
\end{equation}
\cite[(19)]{HHK2}, with the convention that positive torus knots have negative signature. On circle
components HHK fix the absolute grading by declaring one generator to have the grading assigned to the
corresponding generator of Kronheimer and Mrowka's complex~\cite[\S6.1]{HHK2}, \cite[Prop.~4.4]{KM2}.

\begin{theorem}[{\cite[Thm.~5.4]{HHK2}}]\label{thm:hhk54}
Let $L_1\colon S^1\to\Pst$ be a restricted immersed circle that misses the left and right edges of the
pillowcase, and let $\tilde L_1\colon[0,2\pi]\to\R^2$ be a lift of $t\mapsto L_1(e^{it})$. Then its vertical degree
$d=\frac1{2\pi}|\theta(\tilde L_1(2\pi))-\theta(\tilde L_1(0))|$ is even, and $HF(L_0,L_1)$ has rank $d/2$ in each of
the four gradings.
\end{theorem}

\begin{remark}\label{rem:invariance}
By~\cite[Thm.~4.1, Cor.~4.6]{HHK2}, $HF(L_0,L_1)$ is independent of $\epsilon$ and, as a relatively graded group,
depends only on the free homotopy classes of the circle components of $L_1$ and the homotopy class rel
endpoints of its arc. We do not use this: the computations below are direct, and the invariance serves only
as a consistency check in Section~\ref{sec:complex}.
\end{remark}

\subsection{Tangle decompositions}\label{ss:decomp}
Let $K\subset S^3$ be a knot, decomposed along a Conway sphere as $(S^3,K)=(Y,T)\cup(D,U)$, where $(D,U)$ is
the trivial $2$-tangle in a ball and the boundary is identified with $(S^2,\{a,b,c,d\})$~\cite[(17)]{HHK2}.
Let $R(Y,T)$ be the space of conjugacy classes of traceless $\SU$ representations of $\pi_1(Y\setminus T)$
and $R_\pi(Y,T)$ its version perturbed by a holonomy perturbation $\pi$ supported in the interior of $Y$.
Restriction to the boundary gives $L_1\colon R_\pi(Y,T)\to P$. The earring curve $L_0$ of \S\ref{ss:earring}
comes from the trivial side. When $L_1$ is a restricted immersed $1$-manifold and $(L_0,L_1)$ is admissible,
\[
\Hnat(Y,T,\pi)=HF(L_0,L_1),
\]
which by~\cite[p.~34, \S6.1]{HHK2} is the direct sum of the contributions of the components of
$R_\pi(Y,T)$. The two abelian traceless representations of $\pi_1(Y\setminus T)$ are the endpoints of the
arc component; the one that extends over the knot complement maps to the corner $(0,0)$, and its
perturbation gives the generator $r_+$~\cite[\S6.1, \S8]{HHK2}. The remaining generators come in pairs $x^\pm$,
one pair for each transverse intersection point of $L_1$ with the interior of $\Delta$. For a torus knot,
with the tangle side unperturbed and a small perturbation on the side of the earring, these intersection
points correspond to the irreducible traceless representations of the knot group, and there are
$1+|\sigma(K)|$ generators~\cite[Thm.~11.1]{HHK1},~\cite[\S11.10]{HHK2}. For $T(3,n)$, Theorem~\ref{thm:main}
shows that this count persists after a small generic perturbation of the tangle side.
Hedden, Herald and Kirk conjecture:

\begin{conjecture}[{\cite[Conj.~6.5]{HHK2}}]\label{conj:hhk65}
``Given a knot $(X,K)$ in a homology 3-sphere, there exists a 2-tangle decomposition as in Equation (17),
such that for suitably small generic perturbations $\pi$, $L_1:R_\pi(Y,T)\to P$ is a restricted immersed
1-manifold and $H^\natural(Y,T,\pi)$ is isomorphic to the reduced instanton homology $I^\natural(X,K)$.''
\end{conjecture}

The conjecture asserts the existence of one good decomposition. The answer does depend on the
decomposition: different choices of the integers $r,s$ of \S\ref{sec:variety} give different curves
$L_1$~\cite[p.~55]{HHK2}, and with the earring on the other tangle the pillowcase homology of $T(4,5)$ has
the wrong rank~\cite[\S10.3]{CHKK}.

\subsection{S-complexes}\label{ss:scomplex}
Daemi and Scaduto refine the chain complex of reduced singular instanton homology to an
\emph{S-complex}~\cite[Def.~3.21, \S4.1]{DSEquiv}
\[
\tC_*(K)=C_*(K)\oplus C_{*-1}(K)\oplus\Z,\qquad
\td=\begin{pmatrix} d&0&0\\ v&-d&\delta_2\\ \delta_1&0&0\end{pmatrix},
\]
where $C_*(K)$ is generated by irreducible critical points of the singular Chern--Simons functional, the
summand $\Z$ in degree $0$ by the reducible one, $\delta_1$ and $\delta_2$ count instantons to and from the
reducible, and $v\colon C_*\to C_{*-2}$ is defined by holonomy maps on one-dimensional moduli spaces;
degrees are read modulo $4$. The S-structure is $\chi(\alpha,\beta,n)=(0,\alpha,0)$, which satisfies $\chi^2=0$
and $\chi\td+\td\chi=0$. Morphisms of S-complexes are degree-$0$ chain maps commuting with $\chi$ and
inducing the identity on the reducible summand, and S-chain homotopy equivalences are defined
accordingly~\cite[Defs.~3.28,~3.30]{DSEquiv}. The S-chain homotopy type of $\tC(K)$ is an invariant of $K$, its
homology is $\Inat(K)$ up to a shift of grading~\cite[Thm.~3.34, Thm.~8.9]{DSEquiv}, and its local
equivalence class carries the integer invariant $h(K)$; for an S-complex with $d=0$, $\delta_1\neq0$ if and
only if $h>0$~\cite[Prop.~4.15, Rem.~4.2]{DSEquiv}.

For a torus knot $K$ the traceless character variety is finite and nondegenerate, with $|\sigma|/2$ irreducible
points and one abelian point, and the S-complex can be built on it: then $d=0$, $\delta_2=0$ and
$C_*(K)=C_1\oplus C_3$ with
\[
\operatorname{rank}C_1=N_1=\lceil|\sigma|/4\rceil,\qquad \operatorname{rank}C_3=N_3=\lfloor|\sigma|/4\rfloor
\]
\cite[\S1, Cor.~3]{DScaduto},~\cite[\S9.5]{DSEquiv},~\cite{W26}; for $T(3,n)$ the generators, gradings and
differential of $\tC(K)$ are assembled in~\cite{WuebbenI}. In $\tC$ the group $C_1$ occurs in degrees $1$ and
$2$ and $C_3$ in degrees $3$ and $0$. Since $\Inat(K)$ has rank $\|\Delta_K\|_1$, the sum of the absolute values
of the coefficients of the Alexander polynomial~\cite[Prop.~1.4]{KM2},~\cite[Cor.~1.8]{LiYe}, the differential
$\td$ has rank
\[
R=\tfrac12\bigl(1+|\sigma(K)|-\|\Delta_K\|_1\bigr)
\]
over $\Q$, and in the degrees of $\tC(K;\Q)$, with the reducible in degree $0$, the group $\Inat(K;\Q)$ has rank
vector
\begin{equation}\label{eq:gradedInat}
\bigl(N_3+1-\lceil R/2\rceil,\ N_1-\lceil R/2\rceil,\ N_1-\lfloor R/2\rfloor,\ N_3-\lfloor R/2\rfloor\bigr)
\end{equation}
in degrees $0,1,2,3$~\cite{W26}. Kronheimer and Mrowka's absolute grading is obtained by adding $\sigma(K)$
modulo $4$~\cite{W26},~\cite[Thm.~8.9]{DSEquiv}.

For small perturbations the pillowcase complex of a torus knot has $1+|\sigma|$ generators: $r_+$, in grading $\sigma\bmod4$, and a
pair $x_i^-,x_i^+$ with $\gr x_i^+=\gr x_i^-+1$ for each irreducible traceless character, the two points over
the corresponding intersection with $\Delta$. If every bigon runs from some $x_i^-$ either to $r_+$ or to some
$x_j^+$ with $j\neq i$, then
\[
\chi(x_i^-)=x_i^+,\qquad \chi(x_i^+)=\chi(r_+)=0
\]
make the pillowcase complex an S-complex $\pill(K)$ over $\F_2$ with reducible generator $r_+$, in which $x_i^-$
and $x_i^+$ play the roles of the two copies of an irreducible generator~\cite{W26}.

\begin{conjecture}[{\cite{W26}}]\label{conj:pillow}
Let $K=T(p,q)$, and consider the pillowcase complex of~\cite[\S\S10--11]{HHK2} over $\F_2$, for the
decomposition in which the earring lies on the trivial tangle and the integers $r,s$ with $pr+qs=1$ satisfy
$0<r<q$, with zero bounding cochains and a holonomy perturbation small enough that the complex has exactly
the $1+|\sigma|$ generators $r_+$, $x_i^\pm$.
\begin{enumerate}
\item[(a)] Every bigon runs from some $x_i^-$ either to $r_+$ or to some $x_j^+$ with $j\neq i$.
\item[(b)] The S-complex $\pill(K)$, with $\sigma$ subtracted from its grading, is S-chain homotopy
equivalent to $\tC(K;\F_2)$.
\end{enumerate}
\end{conjecture}

The decompositions of~\cite[\S\S10--11]{HHK2} depend on integers $(r,s)$ with $pr+qs=1$. The normalization
$0<r<q$ determines them, since $pr\equiv1\pmod q$, and it is satisfied by all the choices printed
in~\cite[\S11]{HHK2} and by those of Theorem~\ref{thm:main}. Section~\ref{sec:other} shows that it is needed: for
$T(3,4)$ with $(r,s)=(7,-5)$ and zero bounding cochains, the homology has rank $7$ rather than $5$.

Since $h$ is an invariant of the S-chain homotopy type, part~(b) predicts that the component of the pillowcase
differential with values in the span of $r_+$ is nonzero if and only if $\delta_1\otimes\F_2\neq0$~\cite{W26}.
